\chapter{Code}

Code blocks are Pygments-highlighted through minted, wrapped in a box with a
small chrome bar carrying the language and, optionally, a filename. Compile
with \texttt{-shell-escape} or none of this works.

\section{A block with a filename}

\begin{code}{python}{fib.py}
from functools import lru_cache

@lru_cache(maxsize=None)
def fib(n: int) -> int:
    """Nothing clever. The cache does the work."""
    return n if n < 2 else fib(n - 1) + fib(n - 2)

print([fib(n) for n in range(10)])
\end{code}

Pass an empty second argument when the snippet has no file to belong to:

\begin{code}{rust}{}
fn main() {
    let total: u32 = (1..=100).filter(|n| n % 3 == 0 || n % 5 == 0).sum();
    println!("{total}");
}
\end{code}

\section{Other languages}

\begin{code}{haskell}{Primes.hs}
primes :: [Int]
primes = sieve [2..]
  where sieve (p:xs) = p : sieve [x | x <- xs, x `mod` p /= 0]
\end{code}

\begin{code}{sql}{report.sql}
select
  date_trunc('month', created_at) as month,
  count(*) filter (where status = 'shipped') as shipped
from orders
where created_at >= now() - interval '1 year'
group by 1
order by 1;
\end{code}

\section{Terminal sessions}

Shell transcripts get a dark surface, because that is what a terminal looks
like and pretending otherwise helps nobody.

\begin{terminal}
$ latexmk -lualatex -shell-escape main.tex
Latexmk: Run number 1 of rule 'lualatex'
Output written on main.pdf (24 pages, 412918 bytes).
Latexmk: All targets (main.pdf) are up-to-date
\end{terminal}

\section{Unhighlighted text}

For output, logs, or anything Pygments would only get wrong:

\begin{plaincode}
notebook-template/
├── styles/     package files, one concern each
├── fonts/      Inter, Lexend, SF Mono
├── vars/       colours and metadata — edit these
├── sections/   your actual writing
└── main.tex
\end{plaincode}

\section{Inline}

Run \cmd{latexmk -c} to clear the auxiliary files, or \cmd{latexmk -C} to also
remove the PDF. The \texttt{minted} cache lives in \cmd{\_minted-main/} and is
already in \texttt{.gitignore}.

\begin{tip}[Including real files]
\texttt{\textbackslash inputcode\{python\}\{src/model.py\}} pulls a file
straight off disk, so the book never drifts from the code it documents.
\end{tip}

\begin{note}[Changing the theme]
The highlighting style is set once in \texttt{styles/code.sty}
(\texttt{style = friendly}). Run \cmd{pygmentize -L styles} to see the rest —
\texttt{tango}, \texttt{xcode}, and \texttt{perldoc} all print well.
\end{note}
